\newif\ifglobal

%%%%%%%%%%%%%%%%%%%%%%%
%%% Compile to view %%%
%%%%%%%%%%%%%%%%%%%%%%%

%%%%%%%%%%%%%%%%%%%%%%%%%%%%%%%%%%%%%%%%%%%%%%%%%%%
%%% Readme:                                     %%%
%%% https://www.overleaf.com/read/bwdjvfjksvjy/ %%%
%%%%%%%%%%%%%%%%%%%%%%%%%%%%%%%%%%%%%%%%%%%%%%%%%%%

% >>> M?: marks a module setting number or important notice

% >>> M4: Set {./relative-path} to external files for this module <<<
\def \modulefiles {./28-EXTMEMENCR}
% To add an external file, use:
% (Replace '---' with actual filename)
% '\input{\modulefiles/tables/---.tex}' for tables
% '\includegraphics{\modulefiles/figures/---}' for figures
% '\subfile{\modulefiles/---}' for register files


% >>> M1: This module is in global project? <<<
% 'YES' - uncomment; 'NO' - comment out
\globaltrue

\ifglobal
    \documentclass[main\_\_EN.tex]{subfiles}
   
\else
    \documentclass[openany, 10pt]{book}

    % >>> M2: In ./00-shared/preamble-trm-module, also find
    %         settings for visibility of todo notes, watermark <<<
    \usepackage{preamble-trm-module}

    % >>> M3: Temporary labels in English,
    %         you can add your own labels there <<<
    \myexternaldocument{./00-shared/temp-labels-en}

\fi %\ifglobal


\setcounter{secnumdepth}{4}


% >>> M5: If this module needs any updates in
% './00-shared/preamble-trm-module.sty'
% (include additional package, etc.), contact your module PM
% or create the file './00-shared/changelog.tex' make the
% necessary updates yourself and document them in this file <<<


\begin{document}


% Sets language into which pre-defined bits of text are translated
\selectlanguage{english}

% Do not delete this block
\ifglobal
% Add the following front matter if
% this module is outside of global project
\else
    \InputIfFileExists{readme}{}{}
    {\let\clearpage\relax \listoftodos}
    \tableofcontents
    %\listoftables
    %\listoffigures
\fi

%%%%%%%%%%%%%%%%%%%%%%%%%
%%% TRM Module Begins %%%
%%%%%%%%%%%%%%%%%%%%%%%%%



\chapter{External Memory Encryption and Decryption (XTS\_AES)}
\label{mod:extmemencr}
\hypertarget{extmemencr}{}

\section{Overview}
The \chipname{} integrates an External Memory Encryption and Decryption module that complies with the XTS-AES standard algorithm specified in \href{https://ieeexplore.ieee.org/document/4493450}{IEEE Std 1619-2007}, providing security for users' application code and data stored in the external memory (flash). Users can store proprietary firmware and sensitive data (e.g., credentials for gaining access to a private network) to the external flash.

\section{Features}

\begin{itemize}
    \item General XTS-AES algorithm, compliant with IEEE Std 1619-2007
    \item Software-based manual encryption
    \item High-speed auto decryption without software's participation
    \item Encryption and decryption functions jointly enabled/disabled by registers configuration, eFuse parameters, and boot mode
    \item Configurable Anti-DPA
\end{itemize}


\section{Module Structure}
The External Memory Encryption and Decryption module consists of two blocks, namely the Manual Encryption block and Auto Decryption block. The module architecture is shown in Figure \ref{fig:extmemencr-arch-ext-mem-encr-decr}.


\begin{figure}[H]
    \centering
    \includegraphics[width=1\textwidth]{\modulefiles/figures/28-EXTMEMENCR-ext-mem-crypto}
    \caption{Architecture of the External Memory Encryption and Decryption}
    \label{fig:extmemencr-arch-ext-mem-encr-decr}
\end{figure}


The Manual Encryption block can encrypt instructions/data which will then be written to the external flash as ciphertext via SPI1.

In the System Registers (HP\_SYS) peripheral (see \ref{mod:sysreg} \textit{\nameref{mod:sysreg}}), the following three bits in register \hyperref[regdesc:HPSYSTEMEXTERNALDEVICEENCRYPTDECRYPTCONTROLREG]{HP\_SYSTEM\_EXTERNAL\_DEVICE\_ENCRYPT\_DECRYPT\_CONTROL\_REG} are relevant to the external memory encryption and decryption:
\begin{itemize}
    \item \hyperref[fielddesc:HPSYSTEMENABLEDOWNLOADMANUALENCRYPT]{HP\_SYSTEM\_ENABLE\_DOWNLOAD\_MANUAL\_ENCRYPT} 
    \item \hyperref[fielddesc:HPSYSTEMENABLEDOWNLOADG0CBDECRYPT]{HP\_SYSTEM\_ENABLE\_DOWNLOAD\_G0CB\_DECRYPT} 
    \item \hyperref[fielddesc:HPSYSTEMENABLESPIMANUALENCRYPT]{HP\_SYSTEM\_ENABLE\_SPI\_MANUAL\_ENCRYPT} 
\end{itemize}
The XTS\_AES module also fetches two parameters from the peripheral eFuse Controller, which are:\\ \hyperref[fielddesc:EFUSEDISDOWNLOADMANUALENCRYPT]{EFUSE\_DIS\_DOWNLOAD\_MANUAL\_ENCRYPT} and \hyperref[fielddesc:EFUSESPIBOOTCRYPTCNT]{EFUSE\_SPI\_BOOT\_CRYPT\_CNT}. For detailed information, please see Chapter \ref{mod:efuse} \textit{\nameref{mod:efuse}}.


\section{Functional Description}

\subsection{XTS Algorithm} 
\label{sec:extmemncr-xts-algo}

Both manual encryption and auto decryption use the XTS algorithm. During implementation, the XTS algorithm is characterized by a "data unit" of 1024 bits, defined in the Section \textit{XTS-AES encryption procedure} of \\ \href{https://ieeexplore.ieee.org/document/4493450}{\textit{XTS-AES Tweakable Block Cipher}} Standard. For more information about the XTS-AES algorithm, please refer to \href{https://ieeexplore.ieee.org/document/4493450}{IEEE Std 1619-2007}.

\subsection{Key} 
\label{sec:extmemncr-key}
\newcommand{\K}{K\mkern-3mu}

The Manual Encryption block and Auto Decryption block share the same ${\K}ey$ when implementing the XTS algorithm. The ${\K}ey$ is provided by the eFuse hardware and cannot be accessed by users.

The ${\K}ey$ is 256-bit long. The value of the ${\K}ey$ is determined by the content in one eFuse block from BLOCK4 \verb+~+ BLOCK9. For easier description, we define:
\begin{itemize}
\item Block$_A$: the block whose key purpose is EFUSE\_KEY\_PURPOSE\_XTS\_AES\_128\_KEY (please refer to Table \ref{tab:efuse-key-purpose} \textit{\nameref{tab:efuse-key-purpose}}). If Block$_A$ exists, a 256-bit ${\K}ey$\begin{math}_{A}\end{math} is stored in it.
\end{itemize}

There are two possibilities of how the ${\K}ey$ is generated depending on whether Block$_A$ exists or not, as shown in Table \ref{tab:extmemencr-key-possibility}. In each case, the ${\K}ey$ can be uniquely determined by ${\K}ey$\begin{math}_{A}\end{math}.

\begin{longtable}[c]{ |c|c|c|c|c| }
\caption{ ${\K}ey$ generated based on ${\K}ey$\begin{math}_{A}\end{math}}
\label{tab:extmemencr-key-possibility}
\\\hline 
\rowcolor{lightgray}
Block$_A$ & ${\K}ey$  & ${\K}ey$ Length (bit) \\\hline
\endfirsthead
\hline
\rowcolor{lightgray}
Block$_A$ & ${\K}ey$  & ${\K}ey$ Length (bit) \\\hline
\endhead
\endfoot
Yes       &   ${\K}ey$\begin{math}_{A}\end{math}    &   256   \\\hline
No        &   \begin{math}           0^{256}\end{math}    &   256   \\\hline
\end{longtable}
\vspace{-1.5em}
{\bfseries Notes: }\\
``YES" indicates that the block exists; ``NO" indicates that the block does not exist; ``\begin{math}0^{256}\end{math}" indicates a bit string that consists of 256-bit zeros.
Note that using \begin{math}           0^{256}\end{math} as ${\K}ey$ is not secure. We strongly recommend to configure a valid key.


For more information on key purposes, please refer to Table \ref{tab:efuse-key-purpose} \textit{\nameref{tab:efuse-key-purpose}} in Chapter \ref{mod:efuse} \textit{\nameref{mod:efuse}}.

\subsection{Target Memory Space} 
\label{sec:extmemencr-target-space}

The target memory space refers to a continuous address space in the external memory (flash) where the ciphertext is stored. The target memory space can be uniquely determined by two relevant parameters: size and base address, whose definitions are listed below.

\begin{itemize}
    \item Size: the \begin{math} size \end{math} of the target memory space, indicating the number of bytes encrypted in one encryption operation, which supports 16, 32, or 64 bytes.
    \item Base address: the \begin{math} base\_addr\end{math} of the target memory space. It is a 24-bit physical address, with a range of 0x0000\_0000 \verb+~+ 0x00FF\_FFFF. It should be aligned to \begin{math} size \end{math}, i.e., \begin{math} base\_addr \% size == 0 \end{math}.
\end{itemize}

For example, if there are 16 bytes of instruction data that need to be encrypted and written to address 0x130 \verb+~+ 0x13F in the external flash, then the target space is 0x130 \verb+~+ 0x13F, size is 16 (bytes), and the base address is 0x130.

The encryption of any length (must be multiples of 16 bytes) of plaintext instruction/data can be completed separately in multiple operations, and each operation has its individual target memory space and the relevant parameters.

For Auto Decryption blocks, these parameters are automatically determined by hardware. For Manual Encryption blocks, these parameters should be configured by users.

\begin{tiplisting}
\vspace{-2em}
\begin{enumerate}
    The ``tweak" defined in Section \textit{Data units and tweaks} of \href{https://ieeexplore.ieee.org/document/4493450}{IEEE Std 1619-2007} is a 128-bit non-negative integer (\begin{math} tweak \end{math}), which can be generated according to \begin{math} tweak \end{math} = (\begin{math} base\_addr \end{math} \& 0x00FFFF80). The lowest 7 bits and the highest 97 bits in \begin{math} tweak \end{math} are always zero.
\end{enumerate}
\end{tiplisting}

\subsection{Data Writing} 
\label{sec:extmemncr-data-writing}

For Auto Decryption blocks, data writing is automatically applied in hardware. For Manual Encryption blocks, data writing should be applied by users. The Manual Encryption block has a register block which consists of 16 registers, i.e., XTS\_AES\_PLAIN\_\textcolor{red}{\it{n}}\_REG (\textcolor{red}{\it{n}}: 0 \verb+~+ 15), that are dedicated to data writing and can store up to 256 bits of plaintext at a time.

Actually, the Manual Encryption block does not care where the plaintext comes from, but only where the ciphertext will be stored. Because of the strict correspondence between plaintext and ciphertext, in order to better describe how the plaintext is stored in the register block, we assume that the plaintext is stored in the target memory space in the first place and replaced by ciphertext after encryption. Therefore, the following description in this section no longer has the concept of ``plaintext", but uses ``target memory space" instead.

{\bfseries How mapping between target memory space and registers works:}

\newcommand{\f}{\mkern-2mu f\mkern-2mu}
Assume a word in the target memory space is stored in \begin{math} address \end{math}, define $o{\f}{\f}set$ = \begin{math} address \% 64 \end{math}, \textcolor{red}{\it{$n$}} =  $o{\f}{\f}set$/\begin{math}{4}\end{math}, then the word will be stored in register XTS\_AES\_PLAIN\_\textcolor{red}{\it{$n$}}\_REG.

For example, when the \begin{math} size \end{math} is 32, all registers in the register block will be used. The mapping between $o{\f}{\f}set$ and registers now is shown in Table \ref{tab:extmemencr-data-padding}.

\input{\modulefiles/tables/28-EXTMEMENCR-data-padding__EN}

\subsection{Manual Encryption Block}

The Manual Encryption block is a peripheral module. It is equipped with registers and can be accessed by the CPU directly. Registers embedded in this block, the System Registers (HP\_SYS) peripheral, eFuse parameters, and boot mode jointly configure and use this module.

\textbf{ The Manual Encryption block is operational only under certain conditions.} The operating conditions are:

\begin{itemize}

    \item In SPI Boot mode:
    
    If bit \hyperref[fielddesc:HPSYSTEMENABLESPIMANUALENCRYPT]{HP\_SYSTEM\_ENABLE\_SPI\_MANUAL\_ENCRYPT} in register\\ \hyperref[regdesc:HPSYSTEMEXTERNALDEVICEENCRYPTDECRYPTCONTROLREG]{HP\_SYSTEM\_EXTERNAL\_DEVICE\_ENCRYPT\_DECRYPT\_CONTROL\_REG} is 1, the Manual Encryption block can be enabled. Otherwise, it is not operational.
    
    \item In Download Boot mode:
    
    If bit \hyperref[fielddesc:HPSYSTEMENABLEDOWNLOADMANUALENCRYPT]{HP\_SYSTEM\_ENABLE\_DOWNLOAD\_MANUAL\_ENCRYPT} in register\\ \hyperref[regdesc:HPSYSTEMEXTERNALDEVICEENCRYPTDECRYPTCONTROLREG]{HP\_SYSTEM\_EXTERNAL\_DEVICE\_ENCRYPT\_DECRYPT\_CONTROL\_REG} is 1 and the eFuse parameter\\ \hyperref[fielddesc:EFUSEDISDOWNLOADMANUALENCRYPT]{EFUSE\_DIS\_DOWNLOAD\_MANUAL\_ENCRYPT} is 0, the Manual Encryption block can be enabled. Otherwise, it is not operational.
    
\end{itemize}

\begin{tiplisting}
\vspace{-2em}
\begin{enumerate}
Even though the CPU can skip cache and get the encrypted instruction/data directly by reading the external memory, users can by no means access ${\K}ey$.
\end{enumerate}
\end{tiplisting}

\subsection{Auto Decryption Block}

The Auto Decryption block is not a conventional peripheral, so it does not have any registers and cannot be accessed by the CPU directly. The System Registers (HP\_SYS) peripheral, eFuse parameters, and boot mode jointly configure and use this block.

\textbf{The Auto Decryption block is operational only under certain conditions.} The operating conditions are:

\begin{itemize}
    
    \item In SPI Boot mode
    
    If the first bit or the third bit in parameter SPI\_BOOT\_CRYPT\_CNT (3 bits) is set to 1, then the Auto Decryption block can be enabled. Otherwise, it is not operational.
    
    \item In Download Boot mode
    
    If bit HP\_SYSTEM\_ENABLE\_DOWNLOAD\_G0CB\_DECRYPT in register\\ HP\_SYSTEM\_EXTERNAL\_DEVICE\_ENCRYPT\_DECRYPT\_CONTROL\_REG is 1, the Auto Decryption block can be enabled. Otherwise, it is not operational.
    
\end{itemize}

\begin{tiplisting}
\vspace{-2em}
\begin{itemize}

\item When the Auto Decryption block is enabled, it will automatically decrypt the ciphertext if the CPU reads instructions/data from the external memory via cache to retrieve the instructions/data. The entire decryption process does not need software participation and is transparent to the cache. Users can by no means obtain the decryption ${\K}ey$ during the process.

\item When the Auto Decryption block is disabled, it does not have any effect on the contents stored in the external memory, no matter if they are encrypted or not. Therefore, what the CPU reads via cache is the original information stored in the external memory.

\end{itemize}
\end{tiplisting}


\section{ Software Process }

When the Manual Encryption block operates, software needs to be involved in the process. The steps are as follows:

\begin{enumerate}
    
    \item Configure XTS\_AES:
    \begin{itemize}
        \item Set register \hyperref[regdesc:XTSAESPHYSICALADDRESSREG]{XTS\_AES\_PHYSICAL\_ADDRESS\_REG} to \begin{math}base\_addr\end{math}.
        \item Set register \hyperref[regdesc:XTSAESLINESIZEREG]{XTS\_AES\_LINESIZE\_REG} to \begin{math}\frac{size}{32}\end{math}.
    \end{itemize}
    For definitions of \begin{math} base\_addr\end{math} and \begin{math}size\end{math}, please refer to Section \ref{sec:extmemencr-target-space}.
    
    \item Write plaintext instructions/data to the registers block XTS\_AES\_PLAIN\_\textcolor{red}{\it{n}}\_REG (\textcolor{red}{\it{n}}: 0-15). For detailed information, please refer to Section \ref{sec:extmemncr-data-writing}. \\Please write data to registers according to your actual needs, and the unused ones could be set to arbitrary values.
    
    \item Wait for Manual Encryption block to be idle. Poll register \hyperref[regdesc:XTSAESSTATEREG]{XTS\_AES\_STATE\_REG} until it reads 0 that indicates the Manual Encryption block is idle.
    
    \item Trigger manual encryption by writing 1 to register \hyperref[regdesc:XTSAESTRIGGERREG]{XTS\_AES\_TRIGGER\_REG}.
    
    \item Wait for the encryption process completion. Poll register \hyperref[regdesc:XTSAESSTATEREG]{XTS\_AES\_STATE\_REG} until it reads 2. \\ \textit{Step 1 to 5 are the steps of encrypting plaintext instructions/data with the Manual Encryption block using the ${\K}ey$.}
    
    \item Write 1 to register \hyperref[regdesc:XTSAESRELEASEREG]{XTS\_AES\_RELEASE\_REG} to grant SPI1 the access to the encrypted ciphertext. After this, the value of register \hyperref[regdesc:XTSAESSTATEREG]{XTS\_AES\_STATE\_REG} will become 3. 
    
    \item Call SPI1 to write the ciphertext in the external flash (see Section \href{https://docs.espressif.com/projects/esp-idf/en/latest/esp32c6/api-reference/peripherals/spi_flash/index.html#api-reference-flash-encrypt}{API Reference - Flash Encrypt} in \href{https://docs.espressif.com/projects/esp-idf/en/latest/esp32h2/index.html}{ESP-IDF Programming Guide}).
    
    \item Write 1 to register \hyperref[regdesc:XTSAESDESTROYREG]{XTS\_AES\_DESTROY\_REG} to destroy the ciphertext. After this, the value of register \hyperref[regdesc:XTSAESSTATEREG]{XTS\_AES\_STATE\_REG} will become 0.
    
\end{enumerate}

Repeat above steps according to the amount of plaintext instructions/data that need to be encrypted.



\section{Anti-DPA}
\chipname{} XTS\_AES supports Anti-DPA. \\
The XTS-AES algorithm can be divided into two steps, according to \href{https://ieeexplore.ieee.org/document/4493450}{IEEE Std 1619-2007}:
\begin{itemize}
    \item Step 1: Calculating T value. In this section, we define this step as "calculating T".
    \item Step 2: Calculating Cipher/Plain text. In this section, we define this step as "calculating D".
\end{itemize}
Different security levels can be configured through registers:
\begin{itemize}
    \item First we define the below parameters for a better description:
    \begin{itemize}
        \item $select\_reg     = \textrm{\hyperref[fielddesc:XTSAESCRYPTDPASELECTREGISTER]{XTS\_AES\_CRYPT\_DPA\_SELECT\_REGISTER}}$
        \item $reg\_d\_dpa\_en = \textrm{\hyperref[fielddesc:XTSAESCRYPTCALCDDPAEN]{XTS\_AES\_CRYPT\_CALC\_D\_DPA\_EN}}$ 
        \item $efuse\_dpa\_en  = \textrm{\hyperref[fielddesc:EFUSECRYPTDPAENABLE]{EFUSE\_CRYPT\_DPA\_ENABLE}}$
        \item $reg\_anti\_dpa\_level = \textrm{\hyperref[fielddesc:XTSAESCRYPTSECURITYLEVEL]{XTS\_AES\_CRYPT\_SECURITY\_LEVEL}}$
        \item $efuse\_anti\_dpa\_level = \textrm{3}$
    \end{itemize}
    \item Configure the security level of Anti-DPA for the XTS\_AES module:
    $$ Anti\_DPA\_level = select\_reg \ ?\  (reg\_anti\_dpa\_level) : (efuse\_dpa\_en * efuse\_anti\_dpa\_level) $$
    When $Anti\_DPA\_level$ equals to 0, Anti-DPA is disabled. The higher the value of $Anti\_DPA\_level$ is, the stronger the Anti-DPA ability is.
    \item Configure whether or not to enable Anti-DPA when the XTS-AES algorithm is calculating D:
    $$ Anti\_DPA\_enabled\_in\_calc\_D = select\_reg \ ? \ reg\_d\_dpa\_en : efuse\_dpa\_en $$
    If $Anti\_DPA\_level$ is not 0, when $ Anti\_DPA\_enabled\_in\_calc\_D $ equals to 1, Anti-DPA is enabled when XTS-AES algorithm is calculating D.\\
    If $Anti\_DPA\_level$ is not 0, Anti-DPA is always enabled when the XTS-AES algorithm is calculating T.
\end{itemize}
\vspace{-0.5em}
\begin{tiplisting}
Configuring whether or not to enable Anti-DPA will have an impact on the external storage access bandwidth:
\vspace{-0.5em}
\begin{itemize}
    \item When Anti-DPA is enabled during the calculation of D, the read and write bandwidth will be significantly impacted when the Anti-Attack level >= 4.
    \item When Anti-DPA is disabled during the calculation of D, the read and write bandwidth will be significantly impacted when the Anti-Attack level >= 6.
\end{itemize}
\end{tiplisting}



%%%%%%%%%%%%%%%%%%%%%%%%
%%% Register Summary %%%
%%%%%%%%%%%%%%%%%%%%%%%%

\clearpage
\section{Register Summary}
\label{sec:extmemencr-reg-summ}
\hypertarget{extmemencr-reg-summ}{}

The addresses in this section are relative to External Memory Encryption and Decryption base address provided in Table \ref{tab:sysmem-base-address} in Chapter \ref{mod:sysmem} \textit{\nameref{mod:sysmem}}.

The abbreviations given in Column \textbf{Access} are explained in Section \hyperref[glossary-access-types]{\textit{Access Types for Registers}}.

\subfile{\modulefiles/28-EXTMEMENCR-reg-summ__EN}


%%%%%%%%%%%%%%%%%
%%% Registers %%%
%%%%%%%%%%%%%%%%%

\clearpage
\section{Registers}

The addresses in this section are relative to External Memory Encryption and Decryption base address provided in Table \ref{tab:sysmem-base-address} in Chapter \ref{mod:sysmem} \textit{\nameref{mod:sysmem}}.

\subfile{\modulefiles/28-EXTMEMENCR-reg__EN}


\end{document}
